\chapter{Spielerphysik}
\label{ch:players/player-physics}

Die Physik eines Spielers lässt sich mit \emph{Physik-Overrides}
verändern. Damit setzt du Gehgeschwindigkeit, Sprungkraft und
Schwerkraft. Physik-Overrides gelten jeweils für einen einzelnen Spieler
und sind Faktoren: Ein Wert von 2 für die Schwerkraft macht sie zum
Beispiel doppelt so stark.

\section{Ein einfaches Beispiel}\label{ein-einfaches-beispiel}

Hier ein Befehl \texttt{/\allowbreak{}antigravity}, der den aufrufenden
Spieler in geringe Schwerkraft versetzt:

\begin{Shaded}
\begin{Highlighting}[]
\VariableTok{core}\OperatorTok{.}\NormalTok{register\_chatcommand}\OperatorTok{(}\StringTok{"antigravity"}\OperatorTok{,} \OperatorTok{\{}
    \VariableTok{func} \OperatorTok{=} \KeywordTok{function}\OperatorTok{(}\VariableTok{name}\OperatorTok{,} \VariableTok{param}\OperatorTok{)}
        \KeywordTok{local} \VariableTok{player} \OperatorTok{=} \VariableTok{core}\OperatorTok{.}\NormalTok{get\_player\_by\_name}\OperatorTok{(}\VariableTok{name}\OperatorTok{)}
        \ControlFlowTok{if} \KeywordTok{not} \VariableTok{player} \ControlFlowTok{then}
            \ControlFlowTok{return} \KeywordTok{false}\OperatorTok{,} \StringTok{"Nur für Spieler im Spiel"}
        \ControlFlowTok{end}
        \VariableTok{player}\OperatorTok{:}\NormalTok{set\_physics\_override}\OperatorTok{(\{}
            \VariableTok{gravity} \OperatorTok{=} \DecValTok{0.1}\OperatorTok{,} \CommentTok{{-}{-} 10 \% der normalen Schwerkraft}
        \OperatorTok{\})}
    \KeywordTok{end}\OperatorTok{,}
\OperatorTok{\})}
\end{Highlighting}
\end{Shaded}

\section{Verfügbare Overrides}\label{verfuxfcgbare-overrides}

\texttt{player:set\_\allowbreak{}physics\_\allowbreak{}override()}
bekommt eine Tabelle mit Overrides. Laut Lua-API-Referenz\footnote{\url{https://api.luanti.org/class-reference/\#player-only-no-op-for-other-objects}}
sind die wichtigsten:

\begin{itemize}
\tightlist
\item
  \texttt{speed}: Faktor für die Gehgeschwindigkeit (Standard: 1)
\item
  \texttt{jump}: Faktor für die Sprungkraft (Standard: 1)
\item
  \texttt{gravity}: Faktor für die Schwerkraft (Standard: 1)
\item
  \texttt{sneak}: ob der Spieler schleichen kann (Standard:
  \texttt{true})
\end{itemize}

Seit Luanti 5.8 gibt es weitere Overrides, etwa
\texttt{speed\_\allowbreak{}climb}, \texttt{speed\_\allowbreak{}crouch},
\texttt{liquid\_\allowbreak{}fluidity},
\texttt{liquid\_\allowbreak{}sink},
\texttt{acceleration\_\allowbreak{}default} und
\texttt{acceleration\_\allowbreak{}air}. Die vollständige Liste steht in
der Lua-API-Referenz.

\subsection{Altes Bewegungsverhalten}\label{altes-bewegungsverhalten}

Vor Version 0.4.16 gab es den \emph{Sneak-Glitch}, der verschiedene
Bewegungstricks erlaubte. So konnte man mit einer bestimmten Anordnung
von Nodes einen „Aufzug`` hochklettern, indem man schlich (Shift) und
dabei die Leertaste drückte. Das Verhalten war zwar unbeabsichtigt, ist
aber in den Overrides erhalten geblieben, weil viele Server es nutzen.

Um das alte Bewegungsverhalten vollständig wiederherzustellen, braucht
es zwei Overrides:

\begin{itemize}
\tightlist
\item
  \texttt{new\_\allowbreak{}move}: ob der Spieler die neue Bewegung
  verwendet (Standard: \texttt{true})
\item
  \texttt{sneak\_\allowbreak{}glitch}: ob der Spieler
  „Schleich-Aufzüge`` nutzen kann (Standard: \texttt{false})
\end{itemize}

\section{Unverträglichkeit zwischen
Mods}\label{unvertruxe4glichkeit-zwischen-mods}

Achtung: Mods, die denselben Physikwert eines Spielers überschreiben,
vertragen sich meist nicht miteinander. Ein Override überschreibt alle
zuvor gesetzten. Setzen also mehrere Mods die Geschwindigkeit eines
Spielers, gilt nur der zuletzt gesetzte Wert.

\section{Jetzt bist du dran}\label{jetzt-bist-du-dran}

\begin{itemize}
\tightlist
\item
  \textbf{Sonic}: Setze den Geschwindigkeitsfaktor auf einen hohen Wert
  (mindestens 6), wenn ein Spieler dem Spiel beitritt.
\item
  \textbf{Superfeder}: Erhöhe die Sprungkraft, sodass der Spieler 20
  Meter hoch springen kann (1 Meter entspricht 1 Node).
\item
  \textbf{Weltraum}: Lass die Schwerkraft abnehmen, je höher der Spieler
  kommt.
\end{itemize}
